% Subseção da F-RES-OPC (falhas múltiplas e em sequência, RF-RES-05). Dono da sessão:
% Henrique. Fontes: docs/plano-de-desenvolvimento/features/feature-F-RES-OPC.md. Números:
% benchmarks/resultados/falhas-extensoes-2026-11-20.json (comando no JSON).
\subsection{Falhas múltiplas e em sequência}\label{sec:F-RES-OPC}

As falhas são aplicadas uma depois da outra e a rede é medida após cada passo, sempre
contra a rede intacta e com as mesmas origens; ao fim, tudo é reativado na ordem inversa.
As linhas se acumulam: um trecho compartilhado por duas linhas só cai quando as duas
falharam, de modo que o estado final é igual ao de uma falha única com todos os itens
(verificado nos testes). Para $k$ passos o custo é $O((k + 1)\,s\,(V + E)\log V)$. A
página compara a remoção \emph{dirigida} (maior intermediação primeiro) com a
\emph{aleatória} (semente fixa) pela curva da eficiência relativa contra o número de
falhas. No feed real, com 30 origens, 10 remoções dirigidas custam 1,73\% de eficiência,
contra 0,24\% de 10 aleatórias, sem desconectar paradas; cada sequência levou 13,7~s.
\textbf{Limitação:} a intermediação do ataque dirigido é calculada uma vez
na rede intacta; recalculá-la a cada remoção custaria $k$ execuções do Brandes.
